\section{Models}
\label{sec:models}

We adopt an energy accounting model inspired by practices in HPC~\cite{beloglazov2012energy}. Electrical energy consumption is defined as the time integral of instantaneous power consumption,
\[
E = \int_{t_s}^{t_e} P(t)\,dt,
\]
where \(P(t)\) denotes the system power draw and \([t_s, t_e]\) represents the execution interval of a job.

For superconducting quantum computers, system power is dominated by the cryogenic refrigeration and supporting infrastructure required to operate continuously and maintain qubit coherence \cite{matsumoto2024}. However, while this baseline cooling requirement is massive, the total power consumption is not strictly independent of workload characteristics \cite{ueno2025}. We approximated instantaneous power consumption as constant during operation,
\[
P(t) \approx P_{\text{baseline}}.
\]

Under this assumption, the energy consumption of a quantum job simplifies to
\[
E_{\text{job}} \approx P_{\text{baseline}} \cdot (t_e - t_s),
\]
where \(P_{\text{baseline}}\) assumes the average power consumption, including cryogenics, control electronics, and classical infrastructure.

This accounting model simplifies energy attribution without relying on fine-grained gates or pulses, which are difficult to obtain and validate for current quantum hardware. 

The parameters used to model energy consumption in this experiment are summarized in Table~\ref{tab:qcloudsim_energy_params}. Given these parameters, the energy consumption of a quantum job is computed as
\[
E_{\text{job}} = \frac{P_{\text{baseline}} \cdot t_{\text{compute}}}{N_{\text{concurrent}}},
\]
where $N_{\text{concurrent}} = 1$ for exclusive access devices. This formulation accounts for potential time sharing or concurrent execution scenarios without altering the underlying energy model.

Each job consists of a sequence of execution stages that alternate between QPU and CPU devices. For a given job $j$, execution on QPUs and CPUs is recorded as ordered lists of start and finish timestamps:

\[
\{\texttt{qpu\_start}_i^j,\texttt{qpu\_finish}_i^j\}_{i=0}^{N_q-1}
\]
\[
\{\texttt{cpu\_start}_i^j,\texttt{cpu\_finish}_i^j\}_{i=0}^{N_c-1}
\]

where $N_q$ and $N_c$ denote the number of QPU and CPU execution segments, respectively.  

Each QPU segment corresponds to a quantum execution phase, while each CPU segment represents classical pre-processing, post-processing, or intermediate computation required by the hybrid algorithm.

The window $\texttt{qpu\_finish}_i^j - \texttt{qpu\_start}_i^j$ of the $i$-th QPU segment also contains the time the job spends waiting for a connected qubit region (Section~\ref{sec:arch-qpu}); it is retained as the segment's occupancy, while the execution time billed for energy is the circuit execution time recorded by the device,
\[
t^{\text{qpu}}_i = \texttt{qpu\_compute}_i^j,
\]
and the execution time of the $i$-th CPU segment is
\[
t^{\text{cpu}}_i = \texttt{cpu\_finish}_i^j - \texttt{cpu\_start}_i^j.
\]

In that way, quantum and classical execution phases are interleaved multiple times within a single job.

\begin{table}[t]
\centering
\caption{Energy related parameters for superconducting quantum devices}
\label{tab:qcloudsim_energy_params}
\begin{tabular}{llp{5cm}}
\toprule
\textbf{Parameter} & \textbf{Unit} & \textbf{Description} \\
\midrule
$P_{\text{baseline}}$ & kW & Average power consumption, including cryogenic refrigeration, control electronics, and classical infrastructure. \\
$t_{\text{compute}}$ & h & Circuit execution time of a quantum job on the quantum processing unit, excluding time spent waiting for a connected qubit region. \\
$E_{\text{job}}$ & kWh & Energy consumption attributed to an individual quantum job. \\
$N_{\text{concurrent}}$ & -- & Number of concurrently executing jobs sharing the device (if supported). \\
\bottomrule
\end{tabular}
\end{table}

\subsection{Two-Level QPU Allocation}
\label{sec:arch-qpu}

QPUs are the sole spatially constrained resource in the model. Each device loads a topology graph $G=(V,E)$ ($|V|$ physical qubits) and maintains a qubit occupancy map. Admission requires a two-level check: (1) an aggregate counter showing at least $n$ free qubits, and (2) a breadth-first search confirming a \emph{connected} induced subgraph of $n$ free qubits. Upon allocation, selected qubits are marked busy and their adjacent edges cut; upon completion, qubits are freed and cut edges restored between available pairs.

Unlike capacity-only abstractions, this two-level model captures qubit fragmentation: a job meeting the aggregate threshold can still fail admission if free qubits are disconnected, reflecting multi-tenant behavior on sparsely connected superconducting lattices.


QPU service time follows a throughput model based on the vendor-reported Circuit Layer Operations Per Second (CLOPS) metric:

\begin{equation}
\label{eq:qpu-time}
t_{\mathrm{QPU}} \;=\; \alpha\,\frac{M K S \log_2 d}{\mathrm{CLOPS}},
\end{equation}
where $S$ is the shot count, $d$ the circuit depth, $M$ and $K$ the number of circuit templates and parameter updates per submission, and $\alpha$ a unit conversion constant. Per-device calibration tables (readout, single-qubit, and two-qubit gate errors) are additionally loaded to expose an estimated circuit fidelity alongside timing, though fidelity does not feed back into scheduling in the present study. 

\subsection{Classical Device Model}
\label{sec:arch-cpu}

Classical nodes expose two independent capacities, compute units and memory bandwidth, both of which must be satisfied for a job to be admitted; modeling bandwidth separately is necessary because statevector post-processing is frequently bandwidth- rather than compute-bound. The default node derives its service time from the algorithmic cost of the classical half of the variational loop:

\begin{equation}
\label{eq:cpu-time}
t_{\mathrm{CPU}} \;=\; \frac{\kappa\, 2^{\,n} d\, \pi}{\eta\, c^{\,0.85}},
\qquad
\eta = \tfrac{1}{2}\bigl(f_{\mathrm{base}} + f_{\mathrm{boost}}\bigr)\cdot\mathrm{IPC},
\end{equation}

where $\pi$ is the variational parameter count, $\kappa$ a calibration constant, $\eta$ the effective per-core throughput derived from clock and IPC, and the exponent $0.85$ captures sublinear scaling with allocated units. The $2^{\,n}$ term makes the classical side of the loop grow exponentially with qubit width, which is the mechanism by which classical post-processing overtakes quantum execution as circuits widen. A simpler node with randomized service time is also provided for baseline comparisons.

\subsection{Energy and Cost Model}
\label{sec:arch-energy}

Power parameters are set via a single environment configuration with per-device overrides. QPUs draw constant power $P^q$ (dominated by load-independent cryogenic cooling) whenever hosting a job, while classical nodes follow an affine power model based on utilization $u_j$. Instantaneous fleet power is:
\begingroup
\small
\begin{equation}
\label{eq:power}
P(t) \;=\; \sum_{q \in \mathcal{Q}} P^{q}\,\mathbb{1}\!\left[\text{$q$ busy at $t$}\right]
\;+\; \sum_{j \in \mathcal{C}} \Bigl[P^{j}_{\mathrm{idle}} + \bigl(P^{j}_{\mathrm{peak}} - P^{j}_{\mathrm{idle}}\bigr) u_j(t)\Bigr].
\end{equation}
\endgroup

Per-job energy is attributed by phase. For job $J$ with QPU computation times $t^{q}_{i}$ and CPU phase durations $t^{c}_{i}$ on the devices selected in iteration $i$,
\begin{equation}
\label{eq:energy}
E(J) \;=\; \sum_{i=1}^{r}\Bigl(P^{q}_{i} t^{q}_{i} + P^{c}_{i} t^{c}_{i}\Bigr),
\qquad
\mathrm{Cost}(J) \;=\; \rho\, E(J),
\end{equation}
with $\rho$ the electricity price per kWh. Equations~\eqref{eq:power} and~\eqref{eq:energy} are complementary rather than redundant: the former gives the operator's instantaneous load, which determines provisioning and peak demand charges, while the latter gives the per-tenant attribution needed to set prices. Both are driven by the same configuration, so a change to a device's power profile propagates to both views.

\subsection{Telemetry}
\label{sec:arch-telemetry}

Devices stream \texttt{device\_start} and \texttt{device\_finish} events. On each event, the cloud monitor integrates resource-time metrics to record exact utilization and power~\eqref{eq:power}. Exact event-boundary sampling scales with system activity rather than wall-clock time. Concurrently, a records manager logs timestamps to track latency and compute final makespan, energy, and cost~\eqref{eq:energy}. Storing per-iteration traces supports both cross-job and within-job analysis. Aggregated job energy is compared with the integrated power series, which measures fleet occupancy rather than per tenant computation and is not expected to match it.


